% Figure source fragment: taxonomy of efficient attention mechanisms.
% This file is inserted inside a figure environment using \input.
\begin{center}
\setlength{\fboxsep}{5pt}
\renewcommand{\arraystretch}{1.2}
\begin{tabular}{cc}
    \fbox{\parbox[t][1.8cm][t]{0.43\textwidth}{\centering
        \textbf{Hardware-efficient attention}\\[0.2em]
        \scriptsize Same attention function; reduce data movement and
        improve kernel execution.\\
        Example: FlashAttention-2.}} &
    \fbox{\parbox[t][1.8cm][t]{0.43\textwidth}{\centering
        \textbf{Compact attention}\\[0.2em]
        \scriptsize Reduce stored key--value state through sharing or
        compression.\\
        Examples: MQA and GQA.}}\\[0.35em]
    \fbox{\parbox[t][1.8cm][t]{0.43\textwidth}{\centering
        \textbf{Sparse attention}\\[0.2em]
        \scriptsize Select a subset of token or block interactions.\\
        Examples: local, block-sparse, and retrieval-based attention.}} &
    \fbox{\parbox[t][1.8cm][t]{0.43\textwidth}{\centering
        \textbf{Linear attention}\\[0.2em]
        \scriptsize Reformulate similarity computation or use recurrent
        state updates.\\
        Examples: kernelized and gated linear attention.}}
\end{tabular}

\par\smallskip
\footnotesize
Hybrid architectures combine multiple mechanisms across layers or components.
\normalsize
\end{center}
